\section*{The chi-square cap}

The per-variant QC in arm 001 caps chi-square at 80. That single filter turns
out to matter more for this study than any other, and it was found by noticing
that Lp(a)'s \texttt{h2} of $0.007$ is too small even for a tagging problem.

Measured on the published Lp(a) release (9{,}183{,}329 variants with finite
$\beta$ and standard error):

\begin{center}
\begin{tabular}{@{}lr@{}}
\toprule
variants above $\chi^2 = 80$ & 3{,}819 \quad (0.042\%) \\
share of the trait's summed $\chi^2$ they carry & 42.5\% \\
variants in the \emph{LPA} window, chr6:160.0--161.1\,Mb & 3{,}970 \\
\quad of those, removed by the cap & 2{,}894 \quad (72.9\%) \\
largest, rs10455872 & $\chi^2 = 90{,}550$ \\
\bottomrule
\end{tabular}
\end{center}

Four ten-thousandths of the variants carry 42\% of the signal, and the cap
deletes them. rs10455872 exceeds the threshold by a factor of 1{,}132.

\subsection*{Why the cap is there}

It is inherited from LD Score regression, and \texttt{bipred.ldsc.ldsc\_rg}
says so directly: it is a one-step unfiltered LDSC whose regression weights are
built from the fitted means, so a variant whose observed chi-square far exceeds
the LDSC line keeps near-full leverage on the slope. The docstring instructs
the caller to cap extreme chi-square before calling it, and the reference
implementation excludes $\chi^2 > \max(0.001N, 80)$ for the same reason.

The cause is statistical, not numerical. At $N = 284{,}044$ a $\chi^2$ of
$90{,}550$ is a standardised $\hat\beta$ of about $0.56$, well inside the
$|\hat\beta| < 1$ the standardisation requires and nowhere near a precision
limit.

\subsection*{Why it should not apply to the joint fit}

Both drivers build \emph{one} per-variant mask and pass the same variant set to
LD Score regression and to the bivariate fit. So a filter that exists to
protect a linear regression also removes those variants from a Bayesian mixture
model whose slab component exists precisely to hold large effects. bipred does
not need the cap and pays its cost anyway.

\subsection*{It is not confined to Lp(a)}

The obvious guess is that this is an oligogenic-trait problem and that
polygenic traits pay almost nothing. That guess is wrong.
\texttt{../\_lib/chi2\_cap\_audit.py} recomputes the loss for every registered
trait, counting only variants present in the LD reference --- the only ones a
fit could have seen:

\begin{center}
\begin{tabular}{@{}lrrrl@{}}
\toprule
trait & over cap & \% of $\sum\chi^2$ removed & largest $\chi^2$ & at \\
\midrule
bilirubin, total  & 2{,}262 & 73.0 & 107{,}799 & rs887829 \\
bilirubin, direct &   788 & 69.5 &  76{,}937 & rs887829 \\
lipoprotein(a)    &   529 & 50.5 &  12{,}694 & rs3103347 \\
urate             & 1{,}645 & 46.7 &  14{,}649 & rs7675964 \\
cystatin C        & 2{,}039 & 27.0 &  21{,}274 & rs35610040 \\
ALP               & 2{,}960 & 25.5 &  10{,}369 & rs651007 \\
SHBG              & 1{,}957 & 22.9 &   8{,}332 & rs727428 \\
CRP               & 1{,}268 & 17.7 &   5{,}461 & rs4420638 \\
GGT               & 1{,}838 & 16.0 &   7{,}772 & rs5760492 \\
LDL               &   420 &  7.6 &   2{,}013 & rs445925 \\
gout              &   395 &  6.5 &   1{,}907 & rs2231142 \\
VTE               &   153 &  2.2 &     758 & rs651007 \\
gallstones        &   133 &  1.9 &   1{,}402 & rs4953023 \\
CAD               &    29 &  0.4 &     443 & rs2891168 \\
\bottomrule
\end{tabular}
\end{center}

Nine of the fourteen lose more than 15\% of their signal; four lose more than
45\%. Total bilirubin loses 73\%, driven by rs887829 at \emph{UGT1A1} at
$\chi^2 = 107{,}799$. Under 3{,}000 variants out of a million carry that much
of the signal in these traits, and the filter is built to delete exactly them.

The bottom of the table is as informative as the top. CAD loses 0.4\% and LDL
7.6\%, which is why the \texttt{ldl-cad} study --- the one place this pipeline
had been examined closely --- could not have surfaced this. Those two traits
are nearly the least affected of the fourteen.

This has an awkward consequence for the calibration in the first figure.
dbilirubin $\times$ bilirubin is one of the two pairs used to fix what a large
\texttt{rho\_beta} looks like, at $+0.9926$ --- and both of its traits are the
worst affected in the table. The anchor is not damaged in a way that obviously
biases it, since both traits lose the same locus and the surviving effects
still agree, but it is not the clean reference it was presented as.

\subsection*{What it does and does not change}

It does not overturn the ordering in the first figure: \texttt{rho\_beta} still
separates related from unrelated pairs, and it still places Lp(a) $\times$ CAD
with the related ones while \texttt{rg} cannot resolve it. Every pair was
capped identically, so the comparison between pairs is unaffected.

What it does change is the mechanism. The report cannot claim that averaging
one very large locus over a million variants is what defeats \texttt{rg} here,
because that locus was removed before fitting. The defensible claim is
narrower and still worth making: with the dominant locus gone, Lp(a)'s residual
architecture is around 100 causal variants, \texttt{rg} averages that to noise,
and \texttt{rho\_beta} does not. What the uncapped fit would give is unknown,
and is the first thing to run.

\subsection*{A prediction the cap makes about polygenicity}

There is a second consequence, and it points at numbers already in this report.
Truncating the top of the effect-size distribution does not reduce the
heritability the sampler has to account for --- it removes the large effects
that would have accounted for it. The mixture then has to explain the same
signal with small effects, and the only way to do that is to use more of them.
The cap should therefore \emph{inflate} \texttt{n\_causal}.

The estimates are high enough to invite that reading: cystatinC at 21{,}000
causal variants, CRP at 17{,}117--18{,}414, GGT at 13{,}741--14{,}388, out of
roughly $10^6$ fitted.

\textbf{The cross-trait comparison does not support it.} Putting the loss
column beside \texttt{n\_causal}, the three traits that lose most --- bilirubin
at 73\%, Lp(a) at 51\%, urate at 47\% --- come in at roughly 13{,}000, 130 and
3{,}700 causal variants, which is the full range of the table. The two highest
polygenicity estimates, CRP and cystatinC, lose 18\% and 27\%, less than the
median. If the cap were inflating \texttt{n\_causal}, this is where it would
show, and it does not.

That is not a refutation either, because the comparison is confounded: the
traits differ in true polygenicity, sample size and how concentrated their
architecture is, and those differences are larger than the effect being looked
for. The test that would settle it is within-trait --- refit the same trait
with the cap raised and see whether \texttt{n\_causal} moves. Until that is
run, the mechanism is plausible, the cross-trait evidence is against it, and
neither is worth much.
